% annual_transition_extension.tex --- annualised transition variant only

\section{Annual Transition Extension}
\label{sec:annual_transition}

The five-year benchmark remains the paper's main steady-state object, but it compresses the
early family-formation window into a coarse transition block. The annualised branch is useful
because it opens that window back up. It now does so with a cleaner closure than before. On
the upgraded annual permits-starts-completions-stock block, the paper now has a stationary
annual no-drift endpoint plus a full rational-expectations transition around that endpoint.
The annual branch is therefore no longer only a drifting short-horizon stress test. It is a
stationary annual closure with additional short-window support experiments layered on top.

The annual transition still starts from a time-zero cross section that matches the observed
age distribution of four balance-sheet states by construction: renter without debt, renter
with debt, owner with mortgage, and owner outright. The annual cases are then frozen around
two support regimes. The \emph{benchmark} case uses qualification weight $0.25$ and
family-help weight $0.15$. The \emph{robustness} case uses qualification weight $0.40$ and
family-help weight $0.30$. Table~\ref{tab:annual_transition_calibration_targets} collects the
resulting annual transition calibration layer. Those targets still discipline the first five
annual periods most tightly, but they no longer exhaust the annual branch because the same
calibrated construction-flow block now admits a stationary annual endpoint.

Two reduced-form annual entry mechanisms are allowed to operate for young first-time buyers.
The first is a broad help-to-buy / qualification margin that proxies for low-down-payment
products, flexible underwriting, and co-borrower style qualification support. The second is a
narrower family-help margin that acts like a deposit gift at the point of purchase. In the
short-window annual exercises these remain transition objects: they explain how some young
households still enter ownership with mortgage debt despite high house prices, not why house
prices are high in the first place. In the paper's interpretation, high prices still come from
persistent supply restriction; the annual entry channels are coping mechanisms.

\begin{table}[htbp]
\centering
\caption{Annual transition calibration targets and frozen cases}
\label{tab:annual_transition_calibration_targets}
\begin{tabular}{p{3.2cm} p{4.0cm} p{4.6cm} c c c}
\toprule
Group & Moment & Target or role & Drift baseline & Benchmark & Robustness \\
\midrule
Snapshot matched by construction & Owner share 25-34 at t=0 & 0.407 observed t=0 snapshot & 0.407 & 0.407 & 0.407 \\
 & Mortgaged-owner share 25-34 at t=0 & 0.306 observed t=0 snapshot & 0.306 & 0.306 & 0.306 \\
 & Owner share 35-44 at t=0 & 0.605 observed t=0 snapshot & 0.605 & 0.605 & 0.605 \\
 & Mortgaged-owner share 35-44 at t=0 & 0.523 observed t=0 snapshot & 0.523 & 0.523 & 0.523 \\
External incidence targets & Help-to-buy incidence over first 5 years & 0.30-0.40 external first-time buyer help-to-buy band & -- & 0.181 & 0.294 \\
 & Family-help incidence over first 5 years & 0.20-0.30 external first-time buyer family-help band & -- & 0.105 & 0.213 \\
Main annual calibration window & Owner share 25-34 at t=5 & short-run annual transition discipline & 0.610 & 0.657 & 0.684 \\
 & Mortgaged-owner share 25-34 at t=5 & short-run annual transition discipline & 0.367 & 0.407 & 0.430 \\
 & Owner share 35-44 at t=5 & short-run annual transition discipline & 0.739 & 0.770 & 0.789 \\
 & Mortgaged-owner share 35-44 at t=5 & short-run annual transition discipline & 0.315 & 0.338 & 0.351 \\
Persistence check only & Owner share 25-34 at t=20 & persistence check, not main fitting target & 0.680 & 0.718 & 0.740 \\
 & Mortgaged-owner share 25-34 at t=20 & persistence check, not main fitting target & 0.320 & 0.347 & 0.361 \\
 & Owner share 35-44 at t=20 & persistence check, not main fitting target & 0.841 & 0.867 & 0.881 \\
 & Mortgaged-owner share 35-44 at t=20 & persistence check, not main fitting target & 0.196 & 0.205 & 0.209 \\
\bottomrule
\end{tabular}
\end{table}


The stationary no-drift endpoints on the calibrated construction-flow block are now
economically sensible and numerically stable. In the benchmark case the annual stationary
equilibrium has $q^{ss} \approx 1.974$ and young mortgaged ownership for ages 25--34 of about
$0.269$; the corresponding baseline and robustness endpoints are $q^{ss} \approx 1.993$ and
$1.961$. That stationary endpoint matters because it turns the annual branch into a genuine
closure rather than a transition that only lives off an imposed drift path and an eventual
price cap.

Figures~\ref{fig:annual_transition_paths} and~\ref{fig:annual_transition_sensitivity} still
report the drifted support experiments, but they should now be read as short-window policy
diagnostics layered on top of the stationary annual closure rather than as the closure itself.
Under the central $+5$ percent drift calibration, the benchmark regime raises the young owner
share from $0.610$ to $0.657$ at $t=5$ and raises the young mortgaged-owner share from
$0.367$ to $0.407$. The robustness case is stronger, moving those same objects to $0.684$ and
$0.430$. At $t=20$ the effect is smaller but still positive. Across drift rates of $+2$,
$+5$, and $+8$ percent, the benchmark and robustness regimes remain ordered the same way. So
the support margins still matter most in the annual crunch window, even though the live annual
closure is now the stationary no-drift branch.

\begin{figure}[htbp]
\centering
\includegraphics[width=\textwidth]{annual_transition_benchmark_paths.png}
\caption{Annual transition benchmark paths under the central $+5$ percent drift calibration.
These drifted support experiments are now read as short-horizon policy diagnostics layered on
top of the stationary annual closure. The benchmark support regime raises young mortgaged
ownership materially in the main calibration window, while the robustness case provides a
stronger data-facing upper bound.}
\label{fig:annual_transition_paths}
\end{figure}

\begin{figure}[htbp]
\centering
\includegraphics[width=\textwidth]{annual_transition_benchmark_sensitivity.png}
\caption{Annual transition sensitivity across drift calibrations. The benchmark and robustness
regimes remain ordered the same way under $+2$, $+5$, and $+8$ percent annual housing drift.
This bounded grid is now a robustness layer around the stationary annual full-RE closure.}
\label{fig:annual_transition_sensitivity}
\end{figure}

The core annual RE object is now the stationary full-path transition from the observed
snapshot to that endpoint. In the benchmark case the no-drift full-RE path gives
$q_1 \approx 1.001$, $q_5 \approx 1.010$, $q_{10} \approx 1.031$, $q_{20} \approx 1.101$, and
$q_{80} \approx 1.721$. The early and middle years are almost identical to the corresponding
no-RE no-drift anchor, while by $t=80$ the ownership distribution is essentially at its
stationary endpoint and prices are still converging upward toward it. The annualised model can
therefore now be described as a genuine full-RE annual transition around a stationary annual
equilibrium.

At the solver level, the compiled-sidecar Bellman RE workflow also now has a clean staged
$T=4$ diagnostic frontier. The corrected branch-corrector packet survives a clean staged
validation on the intended hybrid path, with the best stage-2 residual at about $0.0039$.
That is a numerical validation benchmark for the transition solver, not yet a claim that the
full long-horizon transition problem has been solved.

For paper-facing exposition, the annual RE objects are best grouped into four layers. First,
the \emph{core closure} is the stationary annual endpoint plus the no-drift full-RE transition
around it. That is the live annual equilibrium object. Second, the \emph{simple bridge
objects} ask what happens when one extra forward-looking margin is added without making the
whole annual system recursive. The fertility-side bridge is the exogenous fertility-demand
shock: relative to no-shock full RE, the benchmark price path is about $0.021$ higher at
$t=5$ and about $0.028$ higher at $t=20$, with only tiny changes in ownership composition. The
policy-side bridge is the imposed regime-path RE exercise. A known switch from baseline
support to benchmark support at $t=5$ leaves the price path almost unchanged relative to the
always-benchmark full-RE case, but it lowers young mortgaged ownership by about $0.026$ at
$t=1$ and about $0.038$ at $t=5$ before the switch arrives.

Third, the \emph{reduced-form recursive extensions} put forecasted aggregate states inside the
fixed point without yet embedding the literal annual household recursion. The main
``beyond-$q$-only'' fertility object is the reduced-form endogenous fertility state. In the
benchmark run the fertility multiplier reaches about $1.026$ at $t=5$ and $1.037$ at $t=20$,
which pushes prices about $0.009$ above the no-shock path at $t=5$ and about $0.055$ above it
at $t=20$, again with only modest ownership effects. On the policy side, the forecasted
political-easing state, the recursive support-state object, and the joint fertility-plus-
politics run all belong in this class: they are genuinely recursive in aggregate states, but
those states are still reduced-form forecast rules rather than full household or voting
problems solved anew each period.

Fourth, there is now a \emph{structural-implied extension}. The structural-fertility operator
version no longer feeds back only through an aggregate birth multiplier. It now also uses
first-birth timing moments to tilt purchase transitions away from ages 25--34 and toward ages
35--44 as prices rise, while higher implied childlessness tilts entrant tenure composition
toward renting. Quantitatively, the enriched structural operator lowers the benchmark annual RE
price path by about $0.026$ at $t=20$, $0.113$ at $t=40$, and $0.349$ at $t=80$ relative to
the no-structural RE benchmark. By $t=80$ the benchmark young purchase scale is down to about
$0.76$, the mid-age purchase scale is up to about $1.15$, and the entrant owner scale is down
to about $0.82$. That is a more genuinely micro structural object than the earlier
birth-multiplier-only bridge. But it is still an operator bridge rather than a literal
household-state recursion, because the Python annual branch does not yet carry explicit parity
or children-at-home states.

The older drifting bounded-RE exercises therefore move down one level in the paper's
hierarchy. The drifting $k=3$ solve remains the main robustness case on the older annual
pressure experiment; the drifting $k=5$ case remains the aggressive robustness version; and
the drifting $k=20$ plus $T=40$ full-path solves are appendix / feasibility material only.
They still show that the upgraded stock-flow block cures the thin-price instability, but the
stationary endpoint plus the $T=80$ full-RE no-drift transition is now the live annual
closure.

The annual branch therefore changes the interpretation of the model rather than overturning the
five-year benchmark. The five-year branch remains the paper's main steady-state equilibrium
object. The annual branch now adds a stationary annual closure, a model-consistent RE
transition around that closure, and a short-horizon housing-access layer that isolates the
crunch years the five-year model smooths over: early family formation, low liquid buffers,
and young ownership entry with mortgage debt. In that sense, the annualised model extends the
paper's main mechanism rather than replacing the baseline five-year result.
